% Public text-only snapshot of the instructor's Module 3 HPO and ablation material.
% Slide upstream: /Users/setup/Library/CloudStorage/Box-Box/Teaching/6050/LaTeX/Module 3-MLP/hpo_slides.tex
% Slide upstream SHA-256: 6deaf680eadf3ed04356207376fa359e5bd8c1fc04c54cd93c1b047676403fba
% Live-session upstream: /Users/setup/Library/CloudStorage/Box-Box/Teaching/6050/3- MLP/live/recording/GMT20250918-002913_Recording.transcript (1).vtt
% Live-session upstream SHA-256: 65a63bd04f7df3c20b21725d4d0b0fa0ed98928686561c45b713a774f568410d
% Local preamble dependencies, TikZ source, executable and product-specific code,
% external assets, private course details, student dialogue, product-current claims,
% unverified numerical demonstrations, and claims without a sufficiently clear
% verification boundary are omitted. The book independently derives its formalism,
% implements its experiment, and checks every reported number.
\documentclass{article}
\usepackage{amsmath}
\usepackage{amssymb}
\usepackage[margin=1in]{geometry}

\title{Hyperparameter Optimization and Ablation Studies}
\author{Heman Shakeri}
\date{}

\begin{document}
\maketitle

\section*{Public Snapshot Boundary}

This text-only snapshot preserves the instructor-authored teaching spine from the
Module 3 deck and its live-session explanation. It is a provenance record rather than
a publishable implementation. Product tutorials, local configuration files, code,
diagrams, external resources, and student-specific discussion are omitted. References
to older Module 3 planning files are not retained because those files contain explicit
D2L-derived prose and code; none of that material crosses this boundary.

The source deck also contains strong edge-of-stability, flatness, and generalization
claims tied to recent tutorial material. Those claims and their schematic figures are
omitted here pending a narrower, primary-source-checked treatment. Likewise, release-
specific tool comparisons, compute percentages, benchmark numbers, and recommended
trial counts are omitted rather than presented as general rules.

\section{From Coding to Experimentation}

In conventional programming, correct code is often expected to produce the specified
output. Deep-learning code can be correct while the learned model performs poorly.
Data preparation, architecture, optimization settings, random variation, and the
training budget all affect the outcome. The lecture therefore frames deep learning as
an empirical pursuit: a practitioner navigates a performance landscape through
controlled, recorded experiments.

Model parameters, including weights and biases, are learned by the inner training
algorithm. Hyperparameters, including learning rate, batch size, regularization, and
schedule, govern that learning process from outside it. The lecture's compact phrase
is ``tuning the trainer, not the model.''

Let $\theta$ denote learned weights and $\lambda$ denote a hyperparameter
configuration. A basic nested objective is

\[
\lambda^*
=\arg\min_{\lambda\in\Lambda}
\mathcal{L}_{\mathrm{val}}\!\left(\theta^*(\lambda)\right),
\qquad
\theta^*(\lambda)
=\arg\min_{\theta}\mathcal{L}_{\mathrm{train}}(\theta;\lambda).
\]

Every evaluation of the outer objective requires an inner training run. The outer
objective is noisy, evaluations are expensive, and the search space can mix
continuous, discrete, and categorical choices.

\section{Ablation and Hyperparameter Optimization}

An ablation poses a research question about a design. It begins with a baseline,
states a hypothesis, removes or modifies a component, and measures the resulting
change. Hyperparameter optimization instead searches for a strong training
configuration within a design.

The live lecture relates tuned design comparison to a job interview: each candidate
should be helped to reach their best state before the final comparison. In that
workflow, architectural settings form the outer comparison and training settings are
optimized inside each setting. The source also repeatedly emphasizes fixed seeds,
fixed data splits, and multiple final seeds so stochastic variation is visible.

The source's rule of thumb is:

\begin{quote}
If a choice is the research question, treat it as the ablation. If it is how to make
that design work well, treat it as hyperparameter optimization.
\end{quote}

This boundary is question-dependent. Batch size, optimizer family, capacity, and even
normalization can be a controlled factor in one study and a tuned configuration in
another.

\section{One Change at a Time, Then Interactions}

Changing one factor at a time supports attribution: if two choices move together, a
raw score difference cannot identify which produced it. The lecture then adds the
important qualification that factors can interact. Removing dropout and removing
normalization together may have an effect different from the sum of their separate
effects. A small factorial design is needed to expose that coupling.

For two binary factors with scores $Y_{00},Y_{10},Y_{01},Y_{11}$, an interaction can
be written as

\[
(Y_{11}-Y_{01})-(Y_{10}-Y_{00}).
\]

A nonzero value means that the effect of the first factor depends on the setting of
the second. The empirical size must still be interpreted against stochastic
variation.

\section{A Budgeted Search Workflow}

The source recommends three broad stages:

\begin{enumerate}
\item perform sanity checks before tuning, including the ability to overfit one small
batch;
\item search broadly with short runs to locate promising regions;
\item narrow the range, spend more on survivors, lock a configuration, and repeat the
final procedure across seeds.
\end{enumerate}

Multiplicative quantities such as learning rate and weight decay are searched on a
logarithmic scale. Grid search is transparent in small spaces. Random search provides
more distinct marginal values per dimension at a fixed trial count. Adaptive methods
use previous trials to decide where to sample next.

Multi-fidelity methods initially allocate small budgets to many configurations, remove
unpromising candidates, and increase the budget for survivors. The pruning rate,
initial budget, maximum budget, and total number of trials remain experimental choices.
They trade broader exploration against the risk of discarding a slow-starting run.

\section{Evaluation Discipline and Reproducibility}

The source distinguishes training, validation, and test roles. Training data fit the
weights. Validation data choose hyperparameters, checkpoints, and designs. The test
set is reserved for the final locked procedure. Repeated optimization of validation
performance can overfit the validation set itself.

Stochastic training makes seeds part of the protocol. The lecture recommends keeping
seeds fixed during controlled searches, then repeating locked configurations across a
panel of seeds for final comparison. Reported spread should be named precisely; a
standard deviation across seeds is not automatically a confidence interval.

Each run should preserve the full resolved configuration, data split, seed, code
revision, budget, learning curves, validation metric, and checkpoint rule. Failed and
pruned trials remain evidence about the search space rather than disposable output.

% [The original product-specific workflow, command lines, YAML listings, dashboards,
% integration code, cluster instructions, and notebook recommendations are omitted.]
% [The original CIFAR-10 MLP code and all claimed 42/48/55 percent results are omitted;
% the book uses a separately pre-tested Fashion-MNIST experiment.]
% [The original GP derivation, acquisition-function details, tool-importance claims,
% automatic-pruning percentages, product defaults, NAS/scaling-law survey, and recent
% tutorial claims are omitted. Primary sources are cited directly in the book.]

\end{document}
